\begin{proof}[Proof of \Cref{lem:p061-rigid-etale-rpl-descent}]
Work on the rigid \(\acute{e}\)tale site \(U_{\mathrm{\acute et}}\) of
de Jong--van der Put.  For an object \(W\to U\), define \(\mathscr P(W)\)
to be the compatible families of \(\mathbf R\)-PL functions on all strict
affinoid G-domains of \(W\).  Thus it is the space of global sections of
the affinoid-local G-sheaf
\[
 A\longmapsto\operatorname{PL}_{\mathbf R}(A^{\mathrm B}).
\]
When \(W\) is affinoid, this is
\(\operatorname{PL}_{\mathbf R}(W^{\mathrm B})\).  This definition covers
all, including globally nongood, objects of the rigid
\(\acute{e}\)tale site.  Pullback makes \(\mathscr P\) a presheaf: on good
strict Berkovich spaces this is
\cite[Lemma~5.4(1)]{GJR26}, and the general rigid object follows by the
affinoid-local definition.

We verify the two conditions of
\cite[Proposition~3.2.2]{dJvdP96}.  First, the restriction of
\(\mathscr P\) to every rigid G-site is a sheaf by construction and
\cite[\S5.2(4)]{GJR26}.  The special rigid G-topology used by
de Jong--van der Put has the same category of sheaves as the usual strong
rigid G-topology by \cite[\S2.1]{dJvdP96}, so there is no change of sheaf
category in this step.

For the second condition, let \(f:W'\to W\) be a surjective finite
\(\acute{e}\)tale morphism of strict affinoid rigid spaces.  We prove
exactness of
\begin{equation}
\label{eq:p061-finite-etale-equalizer}
  0\longrightarrow\mathscr P(W)\longrightarrow\mathscr P(W')
  \mathrel{\substack{\longrightarrow\\[-0.6ex]\longrightarrow}}
  \mathscr P(W'\times_W W').
\end{equation}
For clarity, \(f^{\mathrm B}\) is surjective: finite
\(\acute{e}\)taleness and surjectivity make the affinoid algebra map finite
faithfully flat, so its base change to \(\mathscr H(x)\) is a nonzero finite
algebra for every \(x\in W^{\mathrm B}\), and any maximal-ideal quotient
supplies a point above \(x\).  Injectivity in
\eqref{eq:p061-finite-etale-equalizer} follows.  Now take
\(s'\in\mathscr P(W')\) whose two pullbacks to \(W'\times_WW'\) agree.
We first show that \(s'\) is constant on every Berkovich fibre.  Let
\(y_1,y_2\in (W')^{\mathrm B}\) have common image
\(x\in W^{\mathrm B}\).  Since \(f\) is finite, each
\(\mathscr H(y_i)\) is a finite separable extension of \(\mathscr H(x)\).
Choose a common finite, hence complete, valued extension \(\Omega\) of
\(\mathscr H(x)\) and \(\mathscr H(x)\)-embeddings
\(\iota_i:\mathscr H(y_i)\hookrightarrow\Omega\).  The embeddings are
isometric because the absolute value has a unique extension to a finite
extension of a complete non-Archimedean field.  The bounded homomorphism
\[
 \mathscr H(y_1)\widehat\otimes_{\mathscr H(x)}\mathscr H(y_2)
 \longrightarrow\Omega,\qquad
 a\otimes b\longmapsto\iota_1(a)\iota_2(b),
\]
followed by the absolute value of \(\Omega\), is a bounded multiplicative
seminorm.  Its restriction to each factor is exactly the prescribed absolute
value defining \(y_i\), by the isometry just noted.  It therefore gives a
point \(z\) of
\((W'\times_WW')^{\mathrm B}\) whose two projections are \(y_1\) and
\(y_2\).  Equality of the two pullbacks at \(z\) gives
\(s'(y_1)=s'(y_2)\).

There is therefore a unique function \(s:W^{\mathrm B}\to\mathbf R\)
such that \(s'=s\circ f^{\mathrm B}\).  Both Berkovich affinoid spaces are
compact Hausdorff.  Thus the continuous surjection \(f^{\mathrm B}\) is
closed and is a quotient map.  Since \(s'\) is continuous, the defining
property of the quotient topology shows that \(s\) is continuous.  A
finite \(\acute{e}\)tale morphism is finite, flat, and surjective, so
\cite[Proposition~5.11]{GJR26}, applied with the divisible coefficient
group \(\mathbf R\), shows from
\(s'=f^*s\in\operatorname{PL}_{\mathbf R}(W')\) that
\(s\in\operatorname{PL}_{\mathbf R}(W)\).  This proves the equalizer
condition \eqref{eq:p061-finite-etale-equalizer}.

The criterion \cite[Proposition~3.2.2]{dJvdP96} now says that
\(\mathscr P\) is a sheaf for the rigid \(\acute{e}\)tale topology.  The
section \(g\circ q^{\mathrm B}\) belongs to \(\mathscr P(V)\), and its two
pullbacks to \(V\times_UV\) agree because it is pulled back from the
function \(g\).  Sheaf descent supplies a section \(s\in\mathscr P(U)\)
with \(s\circ q^{\mathrm B}=g\circ q^{\mathrm B}\).  The covering map is
surjective, so \(s=g\).  Hence \(g\) is \(\mathbf R\)-PL.
\end{proof}

\begin{proof}[Proof of \Cref{thm:p061-quasismooth-ph-rpl}]
Fix \(x\in X\).  Goodness and strictness give a strict affinoid
neighbourhood \(U\) of \(x\).  Quasi-smoothness is preserved on analytic
domains.  Under the rigid--Berkovich equivalence recalled in
\cite[\S3.4]{Temkin17}, let \(\mathcal U\) be the rigid space associated to
\(U\); it is smooth, quasicompact, and quasiseparated, and
\(\mathcal U^{\mathrm B}=U\).

Apply \cite[Corollary~3.3.2]{Temkin17}.  There are a finite extension
\(L/K\) and a covering in the rigid \(\acute{e}\)tale topology
\[
                 q:\mathcal V\longrightarrow\mathcal U_L
\]
such that \(\mathcal V\) is affinoid and possesses a strictly semistable
affine formal \(L^\circ\)-model.  This is the precise version needed below.
In particular, we do not replace it by the wording of
\cite[Theorem~3.4.1]{Temkin17}, and we do not assert that \(q\) itself is
finite or flat.

Let \(h_U=h|_U\).  The scalar extension \((h_U)_L\) is classically
pluriharmonic on \(\mathcal U_L^{\mathrm B}\) by
\cite[Proposition~7.2(6)]{GR25Classical}, and hence is continuous by the
opening assertion of that proposition.  Its pullback
\((q^{\mathrm B})^*(h_U)_L\) is classically pluriharmonic by the
arbitrary-morphism
direction of \cite[Proposition~7.2(7)]{GR25Classical}; relative boundary of
the associated quasi-\(\acute{e}\)tale Berkovich morphism causes no problem
for this functorial statement.

The space \(\mathcal V^{\mathrm B}\) is good, quasicompact, and strict and
has the strictly semistable formal model just supplied.  Therefore
\cite[Proposition~7.8]{GR25Classical} gives
\[
                   (q^{\mathrm B})^*(h_U)_L\in
                   \operatorname{PL}_{\mathbf R}(\mathcal V^{\mathrm B}).
\]
The complete \(\acute{e}\)tale descent in
\Cref{lem:p061-rigid-etale-rpl-descent} now gives
\((h_U)_L\in
\operatorname{PL}_{\mathbf R}(\mathcal U_L^{\mathrm B})\).  Since \(L/K\)
is finite, it is already complete and is the completion of an algebraic
extension.  Ground-field descent
\cite[Corollary~5.12]{GJR26}, again with coefficient group \(\mathbf R\),
gives
\(h_U\in\operatorname{PL}_{\mathbf R}(\mathcal U^{\mathrm B})
=\operatorname{PL}_{\mathbf R}(U)\).

Every point \(x\in X\) has such a neighbourhood \(U\), and
\(\mathbf R\)-PL functions form a G-sheaf by
\cite[\S5.2(4)]{GJR26}.  Hence
\(h\in\operatorname{PL}_{\mathbf R}(X)\).

Finally, both \(h\) and \(-h\) are classically psh by definition.  For an
\(\mathbf R\)-PL function on a good strict space,
\cite[Theorem~4.3]{GR25Classical} identifies classical psh with
semipositivity.  Thus both signs of \(h\) are semipositive, which is exactly
harmonicity in \cite[Definition~7.2]{GJR26}.
\end{proof}
